\section{Exact normalized-volume minimization with an lc constraint}
\label{a:minimum}

The principal point of this section is that a constrained minimum is
already an ordinary klt minimum for a sufficiently small positive
perturbation. Convergence to the constraint would not suffice for the
finite-generation argument.

\subsection{The ordinary klt inputs}

For a valuation \(u\) centred at a closed point \(x\) of a
\(p\)-dimensional germ, put
\[
 \mathfrak a_m(u)=\{F:u(F)\geq m\},\qquad
 \vol(u)=\lim_{m\to\infty}
       \frac{p!\,\dim_k(\OO_{U,x}/\mathfrak a_m(u))}{m^p}.
\]
The limit exists for these valuation-ideal families
\cite[Theorem~1.1 and Corollary~5.6]{CutkoskyVolumes}; the local rings
here are excellent and normal, so their completions are reduced.
For a klt pair with discrepancy \(A\), the normalized volume is
\(A(u)^p\vol(u)\).

We use the following established results in their ordinary klt setting.
The minimum is attained at a unique ray, its valuation is quasi-monomial,
and its associated graded algebra is finitely generated and defines a
klt K-semistable log Fano cone. The induced grading again minimizes:
see \cite[Main Theorem and Section~7.3]{BlumExistence},
\cite[Theorem~1.2 and Remark~3.8]{XuQuasimonomial},
\cite[Theorem~1.1]{XuZhuangUniqueness}, and
\cite[Theorems~1.1--1.2]{XuZhuangStable}, with the cone
characterization in \cite[Theorem~3.5]{LiXuStable}.
For the bare klt germ, the estimates
\[
 \widehat{\vol}_U(u)\geq c_U\frac{A_U(u)}{u(\m_x)},\qquad
 u(F)\leq C_U A_U(u)\ord_{\m_x}(F)
\]
are the properness and log Izumi estimates of
\cite[Theorems~1.1--1.2]{LiProperness}.
Their constants may depend on this individual germ.

If the boundary is a rational nonnegative combination of divisors of
regular equations, its specialized boundary is the corresponding
combination of divisors of their initial equations.
The precise simultaneous specialization, including possible component
multiplicities, is proved in Lemma~\ref{a:boundary-specialization}.
Thus all the uses below are of effective rational klt pairs.

\begin{lemma}[A uniform comparison at a minimizer]\label{a:minimizer-comparison}
Let \(b\) minimize normalized volume for a \(p\)-dimensional klt pair.
For every centred quasi-monomial valuation \(u\) and every nonzero
\(F\in\OO_{U,x}\),
\begin{equation}
               u(F)\leq C_p\frac{A(u)}{A(b)}\,b(F),
 \label{eq:A3}
\end{equation}
where \(C_p\) depends only on \(p\). Multiplying the discrepancy function
by a common positive constant does not change this inequality.
\end{lemma}

\begin{proof}
The relative expected-vanishing criterion
\cite[Lemma~3.3, Definition~3.4 and Theorem~3.10]{XuZhuangUniqueness}
gives \(A(u)\geq S(b;u)\). On
\(Q_m=\OO_{U,x}/\mathfrak a_m(b)\), let \(W_m(b;u)\) be the sum of the
jumping weights of the induced \(u\)-filtration, with integer rounding
if desired. Then
\[
 \frac{S(b;u)}{A(b)}
       =\lim_{m\to\infty}\frac{W_m(b;u)}{W_m(b;b)}.
\]
The induced filtration is the image filtration: its dimension at level
\(q\) is
\[
 \dim_k\frac{\mathfrak a_q(u)}
               {\mathfrak a_q(u)\cap\mathfrak a_m(b)}.
\]
Filtering this image by the \(b\)-grading gives the same dimension, so
no compatibility of bases is being assumed.

Suppose \(b(F)>0\), and put
\(k_m=\lfloor m/(2b(F))\rfloor\).
Multiplication by \(F^{k_m}\) embeds
\(\OO_{U,x}/\mathfrak a_{m-k_m b(F)}(b)\) into \(Q_m\).
Every vector in this image has induced \(u\)-level at least \(k_m u(F)\).
The image dimension is asymptotic to \(2^{-p}\dim_k Q_m\), whereas
\(W_m(b;b)\leq m\dim_k Q_m\).
Thus
\[
 \frac{S(b;u)}{A(b)}
       \geq 2^{-p-1}\frac{u(F)}{b(F)}.
\]
The integer-rounding errors vanish in this limit.
This proves \eqref{eq:A3}, for example with \(C_p=2^{p+1}\).
A unit has value zero at both valuations.
\end{proof}

A second proof for canonical-generator germs with Cartier-equation
boundaries appears in Section~\ref{a:comparison-cone-proof}, after the
specialization and torus-field constructions it uses.

\subsection{The constrained problem}

\begin{proposition}[Exact penalization]\label{a:exact-penalization}
Let \(x\in U\) be an affine klt germ of dimension \(p\), and let
\(I_0,J\) be nonzero ideals. Suppose \(c>0\) is rational,
\((U,I_0^c)\) is lc, and an lc quasi-monomial place is centred at \(x\).
Let \(b_0\geq1\) and \(c'\geq0\) be rational, and assume
\[
 a(u)=b_0A_U(u)-c'u(J),\qquad A_U(u)\leq a(u)\leq b_0A_U(u).
\]
The minimum of \(a(u)^p\vol(u)\) among centred quasi-monomial lc places
of \(I_0^c\) is attained at a unique ray. This same ray minimizes
normalized volume for a family of ordinary effective rational klt
pairs for every sufficiently small positive rational parameter.
In particular its associated graded is finitely generated.
It is invariant, after an invariant normalization, under any torus
preserving \(x,I_0,J\).

All smallness thresholds and constants in the proof may depend on the
fixed germ and ideals.
\end{proposition}

\begin{proof}
\emph{From ideals to ordinary klt minimizers.}
Choose a common log resolution of \(U,I_0,J,\m_x\).
Replace the order of each ideal by the average orders of sufficiently
many general elements in a prescribed finite-dimensional generating
space. Write the resulting functions as
\[
                C=A_U-B_0,\qquad a_D=b_0A_U-B_J,
\]
where \(B_0\) and \(B_J\) are rational nonnegative combinations of the
divisors of the chosen regular equations, including the coefficients
\(c\) and \(c'\). The resolving support and these general moving
hypersurfaces form a simple normal crossing divisor \(\Phi\):
choose successively transverse members on its finitely many strata.
Dividing each coefficient among enough members makes every new
moving coefficient sufficiently small.

The old component orders are unchanged. The simple normal crossing
discrepancy tests therefore give
\[
             C\geq0,\qquad a_D\geq A_U/2,\qquad a_D\leq b_0A_U.
\]
They also give \(C=0\) exactly on the original ideal-lc places and
\(a_D=a\) there. Indeed the discrepancy expression for \(C\) has
nonnegative residual term and component terms, strictly positive on
every added moving component, including those used for \(J\).
An ideal-lc place has \(C=0\) by comparison with the fixed ideal orders;
conversely, \(C=0\) forces all these moving orders to vanish and recovers
the original ideal equality. Every chosen equation then has exactly its
ideal's order.

For positive rational \(t'\), sufficiently small, the function
\[
                         \frac{C+t'a_D}{1+t'b_0}
\]
is the discrepancy of the effective rational klt pair with boundary
\[
                    \Delta_{t'}=\frac{B_0+t'B_J}{1+t'b_0}.
\]
Let \(u=u_{t'}\) minimize its normalized volume, with \(u(\m_x)=1\).
Comparison with one fixed centred lc place and \(a_D\geq A_U/2\)
bounds \(A_U(u)^p\vol(u)\) independently of \(t'\).
Properness bounds \(A_U(u)\); log Izumi consequently bounds
\(\vol(u)\) away from zero. The same comparison now gives
\(C(u)=O(t')\).

\smallskip
\noindent\emph{Retraction to the lc strata.}
Let \(\Sigma\) consist of the coefficient-one components for \(C\) on
the fixed resolution. Use the monomial retraction of
\cite[Lemma~4.7 and Corollary~4.8]{JM}: retract \(u\) to the generic
stratum of \(\Sigma\)
through its centre, retaining the same positive normal weights, and
call the retraction \(w\). The component formulas give
\[
 \begin{aligned}
 d_{\log}:=A_{(Y,\Sigma)}(u)&\leq C_1C(u),&
 C(w)&=0,\\
 a_D(w)&\leq a_D(u),&
 w(\m_x)&\geq1-C_1C(u).
 \end{aligned}
\]
To see the first bound, express \(C\) as \(A_{(Y,\Phi)}\) plus the
component values multiplied by their discrepancies. These coefficients
are zero exactly on \(\Sigma\), and positive on the omitted components.
For \(d_{\log}\), the omitted coefficients are instead one.
The formula for \(a_D\) has a nonnegative residual term and positive
component coefficients. The maximal ideal is a fixed component sum.
These observations give all four comparisons. In particular, for small
\(t'\), \(w\) is nontrivial and centred at \(x\).

Retraction has removed the lc defect and has not increased \(a_D\).
To compare the two objectives, we must also control its possible increase
in volume. We prove the following quantitative claim for the minimizers
\(u=u_{t'}\) and their retractions \(w\) constructed above:
\begin{equation}
 u(F)\leq(1+C_2C(u))w(F)
                  \qquad(0\ne F\in\OO_{U,x}).
 \label{eq:A4}
\end{equation}
The constant \(C_2\) is independent of \(F\) and of sufficiently small
\(t'\) on this fixed resolved germ. We first isolate a polynomial's
least normal terms by a differential operator of controlled order, then
bound the value lost under differentiation. A local truncation will give
the claim for every regular \(F\).

\smallskip
\noindent\emph{Polynomial isolation.}
Cover the relevant strata by finitely many etale-coordinate charts
\((x_1,\ldots,x_s,y_1,\ldots,y_{p-s})\), with \(x_i=0\) the components
of \(\Sigma\) through the centre and \(w_i=u(x_i)>0\).
For a polynomial \(G\) of degree at most \(M\) in fixed affine generators
of \(k[U]\) vanishing at \(x\), we claim
\begin{equation}
                      u(G)-w(G)\leq C_3(1+M)d_{\log}.
 \label{eq:A5}
\end{equation}
We will obtain a differential expression with \(u\)-value exactly
\(w(G)\), using \(O(1+M)\) tangential and logarithmic normal
derivatives. The order bound must remain independent of the smallest
positive normal weight; otherwise it would not control the minimizers
as their weights approach a face.

Expand \(G\) normally in the \(x_i\) at the generic stratum.
Etaleness identifies the coefficient field as the finite separable
residue-field lift over \(k(y)\). Its normal Taylor coefficients are
restrictions of regular lifted \(x\)-derivatives divided by factorials;
they are regular on the stratum. Lifted \(y\)-derivatives differentiate
these coefficients without changing normal exponents. At a closed
point this agrees with the expansion in \(x\) and translated tangential
coordinates.

Let \(H\) be the sum of the least-weight normal terms of \(G\).
We first bound every active exponent by \(O(1+M)\), independently of
the smallest positive weight. The extension
\(k(U)/k(x_1,\ldots,x_s,y_1,\ldots,y_{p-s})\) is finite.
Fixed multiplication matrices for the affine generators give a relation
\[
                         \sum_j a_j(x,y)G^j=0
\]
of fixed degree in \(G\), with polynomial coefficient degrees
\(O(1+M)\). Indeed entries in a degree-\(M\) polynomial of these
matrices have denominators bounded by powers of fixed denominators;
the characteristic polynomial and denominator clearing preserve a
linear degree bound.

For a positive normal weight vector \(\lambda\), write
\(f(\lambda)=\min_{\alpha\in\operatorname{Supp}_x(G)}
\langle\alpha,\lambda\rangle\).
At least two different powers \(j,j'\) tie at the least value in the
relation. Choose normal exponents \(\beta,\beta'\) in the finite
supports of \(a_j,a_{j'}\) that attain their respective minimum weights.
Here the supports retain only coefficients nonzero in the stratum field.
The tie gives
\[
 f(\lambda)=
 \frac{\langle\beta-\beta',\lambda\rangle}{j'-j}.
\]
Thus \(f\) takes its values among finitely many linear forms whose
coefficients are \(O(1+M)\): the coefficient degrees bound
\(\beta,\beta'\), and \(|j'-j|\geq1\).
The Taylor minimum is locally finite on the positive orthant, since
the coordinates of \(\lambda\) stay bounded below in a small
neighbourhood and all normal exponents are nonnegative.
Let \(\mathcal A(\lambda)\) be the finite set of active Taylor
exponents. For every coordinate \(i\) and all sufficiently small
\(h>0\), local finiteness gives
\[
 \frac{f(\lambda)-f(\lambda-he_i)}{h}
       =\max_{\alpha\in\mathcal A(\lambda)}\alpha_i.
\]
On this segment, \(f\) agrees with the restriction of one of the
finitely many linear forms above, so the quotient is \(O(1+M)\).
It bounds every active coordinate, including tied exponents.
The permissible \(h\) may shrink as the weights approach a face,
but the bound does not: no division by a normal weight has occurred.

Take the least parts of the relation at the given weights and divide
out the lowest power of \(H\). The resulting polynomial relation in
\(H\) has a nonzero constant term, an \(x\)-polynomial with coefficients
polynomial in \(y\) of degree \(O(1+M)\). At a general closed point of
the centre, not every coefficient of \(H\) can vanish to order exceeding
that bound. Otherwise the relation would force every coefficient of
its nonzero constant term to vanish to that order, impossible for a
nonzero polynomial of the stated degree in etale tangential coordinates.
All other coefficients in the relation are regular there, so their
possible vanishing cannot invalidate this implication.

A tangential derivative of order \(O(1+M)\) therefore makes one active
coefficient a unit. Let its normal exponent be \(\alpha\). For each
normal coordinate apply the product of
\(x_i\partial_{x_i}-j\) over the integers from zero to the active
coordinate bound other than \(\alpha_i\).
This isolates the chosen active monomial, with total differential order
\(O(1+M)\). Both sorts of derivatives preserve normal exponents;
lower-weight coefficients remain zero. To interpret the resulting
congruence in the local ring, work at the chosen closed point of the
centre and put
\[
 I_{>}=\bigl(x^\beta:\textstyle\sum_i\beta_iw_i>w(G)\bigr).
\]
This monomial ideal is finitely generated for the individual positive
weights. In the completion, the output is congruent to \(x^\alpha\)
times a coefficient whose residue on the stratum is a unit. Lift that
residue coefficient from \(\OO_Y/(x_1,\ldots,x_s)\) to an ordinary
local unit \(h\). Since \(x^\alpha(x_1,\ldots,x_s)\subset I_{>}\),
the output minus \(x^\alpha h\) lies in the completed ideal
\(I_{>}\), hence in \(I_{>}\) by faithful flatness. Localize this
congruence at the generic centre of \(u\).
Every ideal generator has \(u\)-value strictly larger than \(w(G)\).
The isolated output consequently has \emph{exactly} \(u\)-value \(w(G)\).
The positive gap here may depend on the valuation; its size is never
used in an estimate.

\smallskip
\noindent\emph{The differential loss.}
Let
\(\lambda_0=\bigwedge_i d\log x_i\wedge\bigwedge_jdy_j\).
For a prime divisor \(E\) centred in the chart,
\(\ord_E(\lambda_0)=A_{(Y,\Sigma)}(E)-1\).
Replacing a factor by \(dG_*/G_*\) multiplies the wedge, up to sign,
by \(x_i\partial_{x_i}G_*/G_*\) or
\(\partial_{y_j}G_*/G_*\). The new wedge has order at least \(-1\):
at a uniformizer of \(E\), each logarithmic differential has only a
possible simple logarithmic pole, and a nonzero wedge can contain that
normal logarithmic factor at most once. Pullbacks of regular tangential
differentials have no poles. Therefore each indicated derivation lowers
value by at most \(A_{(Y,\Sigma)}(E)\), including over a proper
subcentre of the stratum. Approximate real monomial weights on a higher
adapted chart by rational divisorial weights. Continuity of the fixed
function values and discrepancy gives the same inequality for \(u\).
The operator \(x_i\partial_{x_i}-j\) obeys it too, because constants
have value zero and \(d_{\log}\geq0\). Composing the bounded number of
factors loses at most their total order times \(d_{\log}\).
The exact value of the isolated output proves \eqref{eq:A5}.

\smallskip
\noindent\emph{From polynomials to local functions.}
For a nonunit \(F\in\OO_{U,x}\), take
\(M=\lceil3w(F)\rceil+1\).
Modulo \(\m_x^M\), the local ring is spanned by polynomials of degree
less than \(M\) in the fixed affine generators. In particular a unit
denominator is inverted by its finite geometric series in this quotient.
Choose such a polynomial \(G\) representing \(F\).
For small \(t'\), \(w(\m_x)\geq1/2\), so the remainder has
\(u\)-value at least \(M\) and \(w\)-value at least \(M/2>w(F)\).
Thus \(w(G)=w(F)\). Also \(w(F)\geq1/2\), so \(1+M\leq C'w(F)\).
Equation~\eqref{eq:A5} and \(d_{\log}\leq C_1C(u)\) now show that
\[
            u(G)\leq(1+C_2C(u))w(F)<M
\]
for all sufficiently small \(t'\). Hence \(u(F)=u(G)\), proving
\eqref{eq:A4}. Units have value zero at both valuations.

\smallskip
\noindent\emph{Exactness of the minimum.}
We now compare the objectives at \(u\) and \(w\).
Put \(K=1+C_2C(u)\). Equation~\eqref{eq:A4} gives
\(\mathfrak a_m(u)\subseteq\mathfrak a_{m/K}(w)\), and therefore
\(\vol(u)\geq K^{-p}\vol(w)\).
On the fixed germ, properness has bounded
\(a_D(w)\leq b_0A_U(u)\leq M_0\).
Choose \(t'\) so small that \(C_2t'M_0<1\).
If \(C(u)>0\), then
\[
 (C+t'a_D)(u)-K\,t'a_D(w)
     \geq C(u)\bigl(1-C_2t'a_D(w)\bigr)>0.
\]
Together with the volume inequality this contradicts minimality of
\(u\), because \(C(w)=0\). Thus \(C(u)=0\) exactly.

On the constraint, \(a_D=a\), so every sufficiently small parameter
gives a constrained minimum. Conversely any other constrained minimum
would also minimize that ordinary klt objective, whose ray is unique.
The ray is consequently independent of the small parameter.
Stable degeneration gives finite generation. A torus preserving the
ideals and point preserves the constrained objective and its unique ray;
an invariant normalization, for instance by \(u(\m_x)\), makes the
valuation invariant. This proves the proposition.
\end{proof}
